\documentclass[11pt,a4paper]{article}
\usepackage[margin=2.3cm]{geometry}
\usepackage{graphicx}
\usepackage{amsmath}
\usepackage{amssymb}
\usepackage{booktabs}
\usepackage[hidelinks]{hyperref}
\usepackage{microtype}
\title{Complex demography does not explain the ancestry-sorting signal:
shared ancestry, admixture-proportion variation, and space}
\author{}
\date{Module E working note --- \today}

\begin{document}
\maketitle
\vspace{-2em}

\begin{abstract}\noindent
Our recombination-matched neutral null treats the 20 replicate
\textit{F.\ polyctena}\,$\times$\,\textit{F.\ aquilonia} hybrid populations as
independent foundings, and a reviewer rightly asked whether a \emph{richer} neutral
demography could reproduce the empirical signature --- high among-population $F_{ST}$
at ancestry-informative (high-DiagnosticIndex) loci combined with \emph{low} local
linkage disequilibrium. We examine the three natural richer neutral demographies:
(1) a common admixed source that splits into the 20 populations; (2) variation in
admixture proportion among populations; (3) space. \textbf{None can produce the
signal.} Shared ancestry \emph{reduces} differentiation and, by symmetry, cannot
concentrate it at diagnostic loci. Admixture-proportion variation \emph{of the
observed magnitude} is bounded --- analytically and from the data --- to
$\lesssim 5\%$ of the observed diagnostic $F_{ST}$. Space is a milder perturbation
than either and cannot exceed them. The single reason underlying all three is that in
a neutral model the differentiation is \emph{genome-wide} (a low-dimensional ancestry
axis), whereas the empirical differentiation is $\sim\!95\%$ \emph{locus-specific} ---
which is simultaneously why diagnostic $F_{ST}$ is high and why local LD is low. A
model-free \emph{linkage floor} sharpens this: because recombination preserves
between-source allele-frequency covariance, pooled LD cannot fall below its
between-population component, and a genome-wide ancestry axis at the empirical $F_{ST}$
has a floor \emph{above} the empirical total LD --- so no recombinant-source history
built on genome-wide ancestry can reach the data. We also record two methodological
corrections made along the way.
\end{abstract}

\section{The question and three neutral routes}
The empirical data show, relative to the neutral null, \textbf{high differentiation but
low LD} at ancestry-informative loci. A neutral demography could in principle raise
diagnostic $F_{ST}$ three ways: by making the populations share ancestry (route 1), by
letting them differ in overall admixture proportion (route 2), or by arranging them in
space (route 3). We take each in turn. Throughout, ``diagnostic'' loci are
DI\,$>-25$ with parental frequency difference $>0.3$; among-population $F_{ST}$ uses the
20 populations; the neutral simulations are the calibrated recombination burn-in null
(real phased founders, $\mu=0$, empirical cM map, haplodiploid, $K=6250$).

\section{Route 1 --- a common admixed source that splits}
\paragraph{Design and result.} We built a two-stage null: evolve one ancestral admixed
deme for $T_1$ generations, then found 20 sub-demes that drift $T_2 = 50 - T_1$
generations (total $\sim$50, the Nouhaud et al.\ estimate). Sweeping the split time
$T_1 \in \{0,25,45\}$ spans fully independent to near-panmictic. The $F_{ST}$-vs-DI
curve was flat at every $T_1$.

\paragraph{Correction.} That flat result does \emph{not}, on its own, license the
conclusion ``shared ancestry cannot reproduce the signal,'' because a manipulation
check shows the sweep never actually varied realized relatedness. On the full samples,
neutral $F_{ST}$ is flat ($\approx 0.02$) across $T_1$, the among-population coancestry
correlation is pinned at $-1/(n-1)$ (the value for independent demes) at every $T_1$,
and divergence from the ancestral pool is equal for $T_2=25$ and $T_2=5$. The
among-deme differentiation is dominated by the independent founding draw (each sub-deme
redraws its founders), and post-split drift at $K=6250$ over $\le 50$ generations is
negligible; a symmetric split cannot create among-deme covariance structure regardless.
We therefore \textbf{retract} the earlier standalone ``shared-source'' conclusion.

\paragraph{What still holds, structurally.} Shared ancestry \emph{lowers} among-population
differentiation; it cannot create a differentiation \emph{excess} concentrated at
diagnostic loci. And more fundamentally (Route 2), with a single genome-wide admixture
level per population, $F_{ST}$ has no reason to depend on DI at all under \emph{any}
neutral demography.

\section{Route 2 --- variation in admixture proportion (the one that can bite)}
This is the one neutral mechanism that can make $F_{ST}$ depend on DI. If population $i$
has aquilonia ancestry $\alpha_i$, its expected frequency at a locus with parental
difference $\delta = f_{aq}-f_{pol}$ is $p_i = f_{pol} + \alpha_i\,\delta$, so the
among-population variance from admixture alone is $\delta^2\,\mathrm{var}(\alpha)$.
Because $\delta$ is largest at high-DI loci, admixture variation raises diagnostic
$F_{ST}$ specifically. The question is purely quantitative: is the observed spread in
$\alpha$ large enough?

\paragraph{It is not (Fig.~\ref{fig:bound}).} The empirical per-population aquilonia
ancestry spans $[0.35,0.64]$ with $\mathrm{sd}(\alpha)=0.071$. Feeding exactly these
$\alpha_i$ through the pure-admixture model (no locus-specific deviation) predicts a
diagnostic $F_{ST}$ of \textbf{0.015}, against the observed \textbf{0.314} --- admixture
variation accounts for \textbf{4.7\%}. This bound is conservative: sampling noise can
only \emph{inflate} the apparent $\mathrm{sd}(\alpha)$, so the true contribution is if
anything smaller. It converges with two independent measurements: a per-locus regression
of population frequency on $\alpha$ (median $7\%$ variance explained) and the earlier
genome-wide ancestry decomposition ($\sim\!6\%$).

\begin{figure}[h]\centering
\includegraphics[width=0.72\textwidth]{../Figures/moduleE_admixture_bound.pdf}
\caption{Observed among-population $F_{ST}$ vs DiagnosticIndex (orange) against the
ceiling that admixture-proportion variation of the observed magnitude
($\mathrm{sd}\,\alpha = 0.07$) can produce (blue). The neutral admixture ceiling
($0.015$) is $\sim\!20\times$ below the observed diagnostic $F_{ST}$ ($0.314$).}
\label{fig:bound}
\end{figure}

\paragraph{Correction.} An attempt to demonstrate this by simulation --- founding each
deme with a different aquilonia:polyctena \emph{count} ratio --- failed its manipulation
check: realized $\alpha$ converged to $\approx 0.50$ for all demes (uncorrelated with
intent). The reason is model biology: pure-polyctena queens $\times$ pure-aquilonia
males pin every F$_1$ daughter at $50\%$ aquilonia, so the founder \emph{count} sets
diversity, not ancestry. Varying $\alpha$ requires founders of mixed ancestry in both
sexes. We did not pursue that build, because the analytic bound above already settles
the quantitative question the simulation would have asked.

\section{Route 3 --- space}
A spatial stepping-stone with limited migration is a milder perturbation than either a
fully shared source (Route 1) or free variation in admixture proportion (Route 2): it
neither shares more ancestry than a common origin nor spreads $\alpha$ more than the
observed range. It therefore cannot exceed the ceilings above, and we do not simulate it
separately.

\section{Why all three fail is one fact in two spaces}
Every neutral route differentiates populations \emph{genome-wide} --- through a
low-dimensional ancestry/drift axis shared across loci. The data are the opposite:
$\sim\!95\%$ of the diagnostic differentiation is \textbf{locus-specific}. That single
fact explains both halves of the empirical signature at once. High diagnostic $F_{ST}$:
each ancestry-informative locus is differentiated on its own, far beyond the $\sim\!5\%$
a common ancestry axis supplies. Low diagnostic LD: precisely \emph{because} the loci
differentiate independently rather than as shared ancestry blocks, neighbouring
diagnostic markers do not covary --- so local LD stays low. A neutral model gets these
backwards together (low $F_{ST}$, high admixture LD); the data show high $F_{ST}$, low
LD. Locus-specific, divergent selection on ancestry is the mechanism that produces
independent per-locus differentiation without linkage.

\section{The linkage floor: a stricter, model-free bound}
The routes above are ruled out on $F_{ST}$ magnitude. There is a second, independent
bound --- on LD --- that is sharper still, and it disposes of the tempting rescue that
``a long enough recombinant phase could break the LD.'' The rescue fails because
\emph{recombination is frequency-neutral in both directions}: it removes within-source
haplotype LD, but it preserves every source's allele frequency at every locus, and
therefore preserves the \textbf{between-source covariance} of allele frequencies across
adjacent loci. Pooled LD thus has an irreducible floor,
\[ \mathrm{LD_{pooled}(\infty)} \;\approx\; \mathrm{LD_{between}}, \]
that no amount of recombination can go below.

\paragraph{Measuring the floor.} We estimate $\mathrm{LD_{between}}$ by permuting
genotypes among individuals independently at each marker, \emph{within} each population.
This destroys within-population LD while holding every population allele frequency ---
and hence $F_{ST}$ --- fixed; the residual pooled adjacent-marker LD is the floor.
Table~\ref{tab:floor} gives the result for the diagnostic loci. Empirically, observed
adjacent LD is $0.127$ and its floor is $0.043$ (so $66\%$ of the empirical LD is
within-population); the current fixed-$\alpha$ null has a floor of only $0.010$
(its high $0.368$ LD is $97\%$ within-population, i.e.\ removable admixture LD).

\begin{table}[h]\centering
\caption{Diagnostic adjacent-marker LD (median $r^2$): observed vs the between-population
floor recombination cannot remove.}
\label{tab:floor}
\begin{tabular}{lccc}
\toprule
 & observed pooled LD & between-pop floor & within-pop (reducible) \\
\midrule
empirical                       & 0.127 & 0.043 & 0.083 \\
fixed-$\alpha$ null (current)    & 0.368 & 0.010 & 0.358 \\
single ancestry axis @ $F_{ST}{=}0.30$ & --- & \textbf{0.137} & --- \\
\bottomrule
\end{tabular}
\end{table}

\paragraph{The decisive comparison (Fig.~\ref{fig:floor}).} Ask what floor a single
genome-wide ancestry axis would carry \emph{if it reached the empirical diagnostic}
$F_{ST}=0.30$. It requires an $\alpha$-spread $5.9\times$ the observed (already
implausible), and even then its floor is $\mathbf{0.137}$ --- \emph{above} the empirical
observed LD of $0.127$, and more than $3\times$ the empirical floor ($0.043$). A
genome-wide axis makes adjacent loci distinguish the \emph{same} populations too
consistently, so its irreducible floor exceeds the empirical \emph{total} LD: no
recombinant-source model built on genome-wide ancestry can reach the empirical
dissociation. The empirical floor is $3\times$ lower, which means adjacent diagnostic
loci distinguish \emph{different} populations --- the direct signature of locus-specific
differentiation.

\begin{figure}[h]\centering
\includegraphics[width=0.82\textwidth]{../Figures/moduleE_ld_floor.pdf}
\caption{The LD-permutation floor (diagnostic adjacent-marker $r^2$). Recombination can
lower pooled LD only to the between-population floor (blue). A single ancestry axis at
the empirical $F_{ST}$ has a floor ($0.137$) \emph{above} the empirical total LD
($0.127$, dashed); empirical's own floor ($0.043$) is three times lower.}
\label{fig:floor}
\end{figure}

\section{Conclusion and honest caveats}
The three natural richer neutral demographies --- shared common source, admixture-
proportion variation, and space --- are individually insufficient, by a wide and
quantified margin, to reproduce the empirical ancestry-sorting signal. This closes the
``you only tested one demography'' concern for these routes and reinforces the
whole-experiment envelope result (the empirical $F_{ST}$ rise sits $\sim\!100$ standard
deviations outside the neutral envelope).

Two corrections are recorded above: the retraction of the standalone shared-source
conclusion, and the failed admixture-by-count simulation. Both arose from insisting that
a simulation actually produce the variation it claims to test before its output is
interpreted --- a discipline we would keep.

\paragraph{One neutral model remains, and the floor sets its bar.} The linkage floor does
not merely rule things out; it specifies exactly what a surviving neutral model must do:
reach $F_{ST}=0.30$ at diagnostic loci while keeping its between-population floor near the
empirical $0.043$. That excludes not only the three routes above but any
\emph{genome-wide-ancestry} history --- including an ancient, freely-recombining hybrid
\emph{swarm}, whose sources differ in overall ancestry and therefore inherit the high
floor. What survives is narrower: an ancient, \emph{persistently structured} hybrid
metapopulation in which drift and recombination act \emph{concurrently} for long enough
that different source lineages come to carry the high-frequency allele at
\emph{adjacent} loci --- fine-scale, locus-specific source differences, not a shared
ancestry axis. Two stacked demands make this severe: the locus-specific differentiation
requires real generations of drift (which recombination-rate rescaling cannot
accelerate), and keeping those differences finer than the $\sim\!10$\,kb / $0.13$\,cM
marker spacing runs against a $\sim\!523$-generation adjacent-marker LD half-life.
Whether neutral processes can meet both at demographically defensible parameters is the
one genuine open question; the floor ($\le 0.043$ at $F_{ST}=0.30$) is the target any
such simulation must hit.

\section*{Reproducibility}
Analytic bound: \texttt{dev/R/moduleE\_admixture\_bound.R} (Fig.~\ref{fig:bound}) and the
empirical hybrid index \texttt{dev/R/} setup. Shared-source null and its manipulation
check: \texttt{moduleE\_slim/run\_sharedsource.sh}, the model \texttt{OUTPUT\_FOUNDERS}
block, \texttt{dev/R/moduleE\_sharedsource\_relatedness.R}. Admixture-by-count attempt:
\texttt{moduleE\_slim/run\_variable\_alpha.sh}, \texttt{dev/R/moduleE\_varalpha\_setup.R},
\texttt{dev/R/moduleE\_varalpha.R}. Linkage floor (Fig.~\ref{fig:floor}, Table~\ref{tab:floor}):
\texttt{dev/R/moduleE\_ld\_floor.R}. This note supersedes the earlier
\texttt{moduleE\_sharedsource} note, whose standalone conclusion is retracted here.

\end{document}
